% Non inclus dans resume.tex (CV PDF). Projets sélectionnés sur index-fr.html / index-en.html.
% Conserver ce fichier pour référence ; ne pas \input sauf si vous voulez une section Projets dans le PDF.

%-------------------------------------------------------------------------------
% SECTION TITLE
%-------------------------------------------------------------------------------
\cvsection{Projets sélectionnés}

%-------------------------------------------------------------------------------
% CONTENT
%-------------------------------------------------------------------------------
\begin{cventries}

%---------------------------------------------------------
\cventry
    {Outils développeur / Assistants IA de codage} % Project
    {uml-mcp} % Organisation
    {Open Source} % Location
    {2022 -- 2026} % Date(s)
    {
    \begin{cvitems} % Description(s) of project
    \item {Serveur MCP (Model Context Protocol) open source qui transforme des prompts d'architecture ou le contexte d'un dépôt en diagrammes UML et assimilés pour les revues et la documentation technique}
    \item {\textbf{Compétences techniques :} MCP, intégration d'outils LLM, automatisation de documentation technique}
    \end{cvitems}
    }

%---------------------------------------------------------
\cventry
    {Infrastructure / Homelab} % Project
    {Homelab personnel et pratique multi-cloud} % Organisation
    {Auto-hébergé} % Location
    {2020 -- 2026} % Date(s)
    {
    \begin{cvitems} % Description(s) de projet
    \item {Conception et exploitation d'un homelab Proxmox (virtualisation, topologies réseau reproductibles, cycle de vie serveur Dell avec gestion hors bande via iDRAC)}
    \item {Services médias et résidentiels autour de Jellyfin ; interconnexion multi-domicile sécurisée par tunnels WireGuard et maillage Tailscale pour l'administration à distance}
    \item {Automatisation résidentielle avec Home Assistant (capteurs, scénarios et intégrations tierces)}
    \item {\textbf{Compétences techniques :} Proxmox, iDRAC, Jellyfin, WireGuard, Tailscale, Home Assistant, déploiement multi-cloud, supervision et observabilité}
    \end{cvitems}
    }

%---------------------------------------------------------
\cventry
    {Domotique / IoT} % Project
    {hass\_renpho} % Organisation
    {Open Source (fork de neilzilla/hass-renpho)} % Location
    {2022 -- 2026} % Date(s)
    {
    \begin{cvitems} % Description(s) de projet
    \item {Maintient une intégration personnalisée Home Assistant (HACS) qui interroge l'API cloud Renpho (application legacy) et expose poids, composition corporelle, tours de taille et objectifs en capteurs et automatisations}
    \item {Refactor du client API (empreinte client), flux de configuration, proxy HTTP optionnel et explorateur OpenAPI hébergé (\url{https://hass-renpho.vercel.app/docs}) pour valider un compte}
    \item {\textbf{Compétences techniques :} Python, Home Assistant, API REST, HACS, Vercel}
    \end{cvitems}
    }

%---------------------------------------------------------
\cventry
    {IoT / Embarqué} % Project
    {ESP32-7SEG} % Organisation
    {Open Source} % Location
    {2024} % Date(s)
    {
    \begin{cvitems} % Description(s) de projet
    \item {Firmware ESP32 FireBeetle pour afficheur Adafruit 7 segments I2C ; serveur web embarqué pour mettre à jour chiffres et caractères depuis un navigateur sur le réseau local (PlatformIO)}
    \item {\textbf{Compétences techniques :} ESP32, C++, PlatformIO, I2C, serveur web embarqué}
    \end{cvitems}
    }

\end{cventries}
